\documentclass[10pt,a4paper]{amsart}
\usepackage[margin=19mm]{geometry}
\usepackage{amsmath,amssymb,mathtools}
\usepackage[scr=rsfs,cal=euler]{mathalfa}
\usepackage{xfrac}
\usepackage[unicode,hidelinks,bookmarks=false]{hyperref}
\newcommand{\ie}{i.e.\ }
\newcommand{\cf}{cf.\ }
\newcommand{\resp}{respectively\ }
\newcommand{\Subsection}{\subsection}
\providecommand{\nobreakdash}{\nobreak\hbox{-}}
\setlength{\emergencystretch}{2em}
\multlinegap0pt
\usepackage[shortlabels]{enumitem}
\usepackage{tikz}
\usepackage{xcolor}
\usepackage{xfrac}
\usepackage{amssymb}
\usepackage{mathtools}
\usepackage[normalem]{ulem}
\newtheorem{proposition}{Proposition}
\newtheorem{lemma}{Lemma}
\newcommand{\A}{\mathbb A}
\newcommand{\R}{\mathbb R}
\newcommand{\T}{\mathbb T}
\newcommand{\Z}{\mathbb Z}
\newcommand{\ep}{\varepsilon}
\newcommand{\norm}[1]{ \Vert #1 \Vert}
\title{Non-unique solutions to the periodic \textsf{gKdV} equation}
\author{Nicholas Gismondi, Kunyi (Mark) Ma, Mandon Pathak, Alexandru F. Radu}
\date{Source: arXiv:2606.06916v1 (CC BY 4.0)}
\begin{document}
\maketitle

\noindent\emph{Source and licence.} Non-unique solutions to the periodic \textsf{gKdV} equation. Version arXiv:2606.06916v1 (CC BY 4.0), distributed by arXiv under the Creative Commons Attribution 4.0 International licence, http://creativecommons.org/licenses/by/4.0/, as recorded in the arXiv licence metadata for this exact version and shown on the abstract page https://arxiv.org/abs/2606.06916v1. Full licence notice: \texttt{licenses/arxiv-2606.06916-CC-BY-4.0.txt}.\par\smallskip

\noindent\textit{Editorial scope.} Counted unit: the article's lemma lem:horrible1 (source0.tex lines 2477--2508). The article's complete original proof of it is delivered with it (source0.tex lines 2509--2838, one \texttt{proof} environment of 330 lines). No mathematics is written, repaired or re-derived: the statement and the proof are the article's own lines, in the article's own notation, and no retained line is substituted or re-flowed.  Retained uncounted as the source-local context the proof needs: lines 860--862; lines 942--985.  Nothing was omitted from any retained span, so the package carries no omission record.  No retained line contains a citation, so the wrapper carries no bibliography and no bibliography entry.  No retained line contains a figure environment, an image-inclusion command or drawing code, so no image is delivered and the package declares no asset dependency.  Editorial formatting only, in addition: an AMS \texttt{amsart} preamble that restates the article's own declarations and macros used by the retained lines, namely the environments \texttt{proposition}, \texttt{lemma} on separate counters, together with the standard packages those lines need.  The retained environments print in this document as Proposition 1, Lemma 1 in order of appearance, whereas the article prints the same items with the numbering of its own full document.  The complete native source of the article as distributed in the arXiv e-print for this exact version is bundled unchanged as \texttt{sources/202220/source0.tex}.
\begin{equation}\label{eq:relaxed}
    \partial_t u + \partial_{xxx} u + \sigma \partial_x\left(u^k\right) = \partial_x E,
\end{equation}

\begin{proposition}\label{prop:ind}
    Fix $0 < \beta < \frac{1}{16}\left(1 - \frac{1}{k}\right)$. There are sequences $u_q:[0,1] \times \T \to \R$, $E_q : [0,1] \times \T \to \R$, and $\lambda_q \in \R$, $q \geq 0$, such that the following hold:
    \begin{enumerate}
        \item\label{i:1} Each $\lambda_q \geq 1$ is a power of $2$, $\lambda_{q'-1} < \lambda_{q'}$ for all $q' \leq q$ ($\lambda_{-1}=0$), and
        $$
        \sum_{q' \leq q}\lambda_{q'}^{-\beta} \leq \frac{1}{4} - 2^{-q-100}.
        $$
        \item\label{i:2} $(u_q,E_q)$ is a smooth solution to the relaxed equation given by ~\eqref{eq:relaxed}. In addition, $u_q$ has zero mean and the temporal support of both $u_q$ and $E_q$ is contained in
        $$
        \left(\frac{1}{2}- \sum_{q'\leq q} \frac{1}{\lambda_{q'}^\beta},\frac{1}{2} + \sum_{q' \leq q} \frac{1}{\lambda_{q'}^\beta}\right);
        $$
        \item\label{i:3} Upon setting $u_{-1}=0$, for every $q' \leq q$, the frequency support of $u_{q'} - u_{q'-1}$ lies in the annulus
        $$
        \left\{\xi : 2\lambda_{q'-1}^{3/2} < |\xi| < \lambda_{q'}^{3/2}\right\}.
        $$
        Also for every $1\le p<k$ and every
$q'\le q$ we have
\[
    \|u_{q'}-u_{q'-1}\|_{C_t^0L_x^p}
    \lesssim_p
    2^{-C(p)q'}.
\]
When $k\ge3$, for every $0<\alpha<\frac12-\frac1k$ and every $q'\le q$ we also
have
\[
    \|u_{q'}-u_{q'-1}\|_{C_t^0\dot H_x^\alpha}
    \lesssim_\alpha
    2^{-C(\alpha)q'}.
\]
In addition,
\[
    \|u_q\|_{C_t^0L_x^1}>\delta(1+2^{-q})
\]
for some $\delta>0$ chosen independent of $q$.
\item\label{i:4} The error $E_q$ satisfies the estimate
        \[
        \norm{E_q}_{\dot{\mathbb{A}}_{x,t}^{-s}} < 2^{-q-200}.
        \]
        \item\label{i:5} There is $C>0$ such that
        $$
        \sum_{\xi_1 + \ldots + \xi_k \not = 0} \left\Vert  \hat{u}_{q}(t,\xi_1) \hat{u}_{q}(t,\xi_2) \cdots \hat{u}_{q}(t,\xi_k) \right\Vert_{C_t^0}|\xi_1 + \ldots + \xi_k|^{-s} < C - 2^{-q-100}.
        $$
    \end{enumerate}
\end{proposition}

\begin{lemma}\label{lem:horrible1}
    We have that
    \[
\begin{aligned}
AF(w_{q+1},j)
\lesssim_{q,j,s}\;&
\lambda_{q+1}^{(1-\epsilon_{q+1})(\frac jk-1)-s\epsilon_{q+1}}
A_{q+1}^{j/k}
\\
&+
\lambda_{q+1}^{(1-\epsilon_{q+1})(\frac jk-1)}
A_{q+1}^{\frac jk-1}
\|E_q\|_{\dot{\mathbb \A}^{-s}_{x, t}}
\\
&+
\lambda_{q+1}^{(1-\epsilon_{q+1})(\frac jk-1)-s\epsilon_{q+1}}
A_{q+1}^{\frac jk-1}
\\
&+
\lambda_{q+1}^{(1-\epsilon_{q+1})(\frac jk-1)}
A_{q+1}^{\frac jk-2}.
\end{aligned}
\]
In particular,
\[
AF(w_{q+1},j)
\lesssim_{q,j,s}
\lambda_{q+1}^{(1-\epsilon_{q+1})(\frac jk-1)}
A_{q+1}^{j/k}.
\]

\end{lemma}

\begin{proof}
    Clearly
    \begin{equation*}
        \begin{split}
            &\Vert \hat{w}_{q+1}(t,\eta_1) \cdots \hat{w}_{q+1}(t,\eta_j) \Vert_{C_t^0}\\
            &\lesssim \sum_{n_1,\ldots,n_j = 0}^\infty \sum_{\xi_1, \ldots,\xi_j \in \Z} \lambda_{q+1}^{j(1-\epsilon_{q+1})\left(\frac{1}{k}-1\right)} \left|\hat{\phi}_k\left(\lambda_{q+1}^{\epsilon_{q+1}-1}\xi_1\right) \cdots \hat{\phi}_k\left(\lambda_{q+1}^{\epsilon_{q+1}-1}\xi_j\right)\right| \left|\binom{1/k}{n_1} \cdots \binom{1/k}{n_j}\right| A_{q+1}^{\frac{j}{k}-n_1 -\ldots - n_j}\\
            &\times \left\Vert \left(E_q^{n_1} h_{q+1}\right)^{\wedge}(t,\eta_1 - \lambda_{q+1}^{\epsilon_{q+1}} \xi_1) \right\Vert_{C_t^0} \cdots \left\Vert \left(E_q^{n_1} h_{q+1}\right)^{\wedge}(t,\eta_1 - \lambda_{q+1}^{\epsilon_{q+1}} \xi_j) \right\Vert_{C_t^0}
        \end{split}
    \end{equation*}
    There are four cases we consider.

    \textbf{Case 1:} $n_1 = n_2 = \ldots =n_j = 0$.\\
    In this case we have that
    $$
    \left(E_q^0 h_{q+1}\right)^{\wedge}(t,\eta_i - \lambda_{q+1}^{\epsilon_{q+1}} \xi_i) = \delta(\eta_i - \lambda_{q+1}^{\epsilon_{q+1}} \xi_i) h_{q+1}(t)
    $$
    and thus
    $$
    \left\Vert \left(E_q^{0} h_{q+1}\right)^{\wedge}(t,\eta_i - \lambda_{q+1}^{\epsilon_{q+1}} \xi_i) \right\Vert_{C_t^0} = \left\Vert \delta(\eta_i - \lambda_{q+1}^{\epsilon_{q+1}} \xi_i) h_{q+1}(t) \right\Vert_{C_t^0} = \delta(\eta_i - \lambda_{q+1}^{\epsilon_{q+1}} \xi_i)
    $$
    for all $1 \leq i \leq j$. Hence
    \begin{equation}\label{eq:n=0_case}
        \begin{split}
            &\sum_{\eta_1 + \ldots + \eta_j \not=0} \sum_{\xi_1,\ldots,\xi_j \in \Z} \lambda_{q+1}^{j(1-\epsilon_{q+1})\left(\frac{1}{k} - 1\right)} \left|\hat{\phi}_k\left(\lambda_{q+1}^{\epsilon_{q+1}-1}\xi_1\right) \cdots \hat{\phi}_k\left(\lambda_{q+1}^{\epsilon_{q+1}-1}\xi_j\right)\right| A_{q+1}^{j/k}\\
            &\times \delta(\eta_1 - \lambda_{q+1}^{\epsilon_{q+1}} \xi_1) \cdots \delta(\eta_j - \lambda_{q+1}^{\epsilon_{q+1}} \xi_j) |\eta_1 + \ldots + \eta_j|^{-s}\\
            &= \sum_{\xi_1 + \ldots + \xi_j \not=0} \lambda_{q+1}^{j(1-\epsilon_{q+1})\left(\frac{1}{k}-1\right) - s\epsilon_{q+1}} \left| \hat{\phi}_k\left(\lambda_{q+1}^{\epsilon_{q+1}-1}\xi_1\right) \cdots \hat{\phi}_k\left(\lambda_{q+1}^{\epsilon_{q+1}-1}\xi_j\right)\right| A_{q+1}^{j/k} |\xi_1 + \ldots + \xi_j|^{-s}\\
            &= \lambda_{q+1}^{(1-\epsilon_{q+1})\left(\frac{j}{k}-1\right)-s\epsilon_{q+1}} A_{q+1}^{j/k} \sum_{\Xi \not=0} |\Xi|^{-s} \sum_{\xi_1,\ldots,\xi_{j-1}} \left|\prod_{i=1}^{j-1}\hat{\phi}_k\left(\lambda_{q+1}^{\epsilon_{q+1}-1}\xi_i\right) \hat{\phi}_k\left(\lambda_{q+1}^{\epsilon_{q+1}-1}\left(\Xi - \xi_1 - \ldots - \xi_{j-1}\right)\right)\right|\\
            &\times \lambda_{q+1}^{(j-1)(\epsilon_{q+1}-1)}\\
            &\lesssim \lambda_{q+1}^{(1-\epsilon_{q+1})\left(\frac{j}{k}-1\right)-s\epsilon_{q+1}} A_{q+1}^{j/k} \sum_{\Xi \not=0} |\Xi|^{-s} \int_{\R^{j-1}} \left| \prod_{i=1}^{j-1} \hat{\phi}_k(\xi_i) \hat{\phi}_k(\lambda_{q+1}^{\ep -1}\Xi - \xi_1 - \ldots - \xi_{j-1})\right|\, d\xi_1\, \ldots\, d\xi_{j-1}\\
            &\lesssim \lambda_{q+1}^{(1-
            \epsilon_{q+1})\left(\frac{j}{k}-1\right)-s\epsilon_{q+1}} A_{q+1}^{j/k}
        \end{split}
    \end{equation}
    where we have utilized that $s > 1$, our choice of taking $\lambda_{q+1}$ arbitrarily large, and the generalized Young's convolution inequality to get that
    $$
    \left\Vert \left(|\hat{\phi}_k| \ast |\hat{\phi}_k| \ast \ldots \ast |\hat{\phi}_k|\right)\left(\lambda_{q+1}^{\epsilon_{q+1}-1} \cdot\right)\right\Vert_{L^\infty} \lesssim \left\Vert \hat{\phi}_k\right\Vert_{L^1}^{j-1} \left\Vert \hat{\phi}_k \right\Vert_{L^\infty} \lesssim 1.
    $$

    \textbf{Case 2:} $n_i = 1$ for some $1 \leq i \leq j$, $n_m=0$ for all $m \not=i$, $\xi_1 + \ldots + \xi_j = 0$.

    By symmetry considerations let us suppose that $i = 1$; i.e. $n_1 = 1$ and $n_2=\ldots =n_j =0$. Then we have that
    $$
    \left\Vert \left(E_qh_{q+1}\right)^{\wedge}(t,\eta_1 - \lambda_{q+1}^{\epsilon_{q+1}} \xi_1)\right\Vert_{C_t^0} = \left\Vert \hat{E}_q(t,\eta_1 - \lambda_{q+1}^{\epsilon_{q+1}} \xi_1)\right\Vert_{C_t^0}
    $$
   as well as
   $$
   \left\Vert \left(E_q^{0} h_{q+1}\right)^{\wedge}(t,\eta_i - \lambda_{q+1}^{\epsilon_{q+1}} \xi_i) \right\Vert_{C_t^0} = \delta(\eta_i - \lambda_{q+1}^{\epsilon_{q+1}} \xi_i)
   $$
   for all $1 < i \leq j$. Hence, using the same convolution estimate from Case 1, we have
   \begin{equation}\label{eq:n=1_case}
       \begin{split}
           &\sum_{\eta_1 + \ldots + \eta_j \not=0} \sum_{\xi_1 + \ldots + \xi_j=0} \lambda_{q+1}^{j(1-\epsilon_{q+1})\left(\frac{1}{k} - 1\right)} \left|\hat{\phi}_k\left(\lambda_{q+1}^{\epsilon_{q+1}-1}\xi_1\right) \cdots \hat{\phi}_k\left(\lambda_{q+1}^{\epsilon_{q+1}-1}\xi_j\right)\right| A_{q+1}^{\frac{j}{k}-1}\\
           &\times \left\Vert \hat{E}_q(t,\eta_1 - \lambda_{q+1}^{\epsilon_{q+1}} \xi_1)\right\Vert_{C_t^0}\delta(\eta_2 - \lambda_{q+1}^{\epsilon_{q+1}} \xi_2) \cdots \delta(\eta_j - \lambda_{q+1}^{\epsilon_{q+1}} \xi_j) |\eta_1 + \ldots + \eta_j|^{-s}\\
           &= \sum_{\xi_1 + \ldots + \xi_j = 0}  \sum_{\eta_1 \not= \lambda_{q+1}^{\epsilon_{q+1}} \xi_1} \lambda_{q+1}^{j(1-\epsilon_{q+1})\left(\frac{1}{k}-1\right)}A_{q+1}^{\frac{j}{k}-1} \left|\hat{\phi}_k\left(\lambda_{q+1}^{\epsilon_{q+1}-1}\xi_1\right) \cdots \hat{\phi}_k\left(\lambda_{q+1}^{\epsilon_{q+1}-1}\xi_j\right)\right|\\
           &\times \left\Vert \hat{E}_q(t,\eta_1 - \lambda_{q+1}^{\epsilon_{q+1}} \xi_1)\right\Vert_{C_t^0} \left|\eta_1 -\lambda_{q+1}^{\epsilon_{q+1}} \xi_1\right|^{-s}\\
           &= \sum_{\xi_1 + \ldots + \xi_j = 0} \lambda_{q+1}^{j(1-\epsilon_{q+1})\left(\frac{1}{k}-1\right)} A_{q+1}^{\frac{j}{k}-1} \left|\hat{\phi}_k\left(\lambda_{q+1}^{\epsilon_{q+1}-1}\xi_1\right) \cdots \hat{\phi}_k\left(\lambda_{q+1}^{\epsilon_{q+1}-1}\xi_j\right)\right| \Vert E_q \Vert_{\dot{\mathbb{A}}_{x,t}^{-s}}\\
           &\lesssim \lambda_{q+1}^{(1-\epsilon_{q+1})\left(\frac{j}{k}-1\right)} A_{q+1}^{\frac{j}{k}-1} \Vert E_q \Vert_{\dot{\mathbb{A}}_{x,t}^{-s}}.
       \end{split}
   \end{equation}
    
    \textbf{Case 3:} $n_i = 1$ for some $1 \leq i \leq j$, $n_m=0$ for all $m \not=i$, $\xi_1 + \ldots + \xi_j \not= 0$.

    Again from symmetry considerations we assume that $n_1 = 1$ and $n_2 = \ldots = n_j = 0$. From item \ref{i:3} and the fact $(u_q,E_q)$ solve \eqref{eq:relaxed}, the frequency support of $E_q$ is contained in $B(0,k\lambda_q^{3/2})$. Hence we have that
    $$
    |\eta_1 - \lambda_{q+1}^{\epsilon_{q+1}} \xi_1| \lesssim \lambda_q^{3/2}.
    $$
    Since $\eta_i = \lambda_{q+1}^{\epsilon_{q+1}} \xi_i$ for all $2 \leq i \leq j$ then we see that we may choose $\lambda_{q+1}$ large enough to ensure that
    \begin{equation}\label{eq:AF-case3-gap}
        \begin{split}
            |\eta_1 + \eta_2 + \ldots + \eta_j| &= \left|\eta_1 - \lambda_{q+1}^{\epsilon_{q+1}} \xi_1 + \lambda_{q+1}^{\epsilon_{q+1}}\left(\xi_1 + \xi_2 + \ldots + \xi_j\right)\right|\\
            &\geq \lambda_{q+1}^{\epsilon_{q+1}} \left|\xi_1 + \ldots + \xi_j\right| - |\eta_1 - \lambda_{q+1}^{\epsilon_{q+1}} \xi_1|\\
            &\geq \frac{1}{2}\lambda_{q+1}^{\epsilon_{q+1}} \left|\xi_1 + \ldots + \xi_j\right|\\
            &\simeq \lambda_{q+1}^{\epsilon_{q+1}} \left|\xi_1 + \ldots + \xi_j\right|
        \end{split}
    \end{equation}
    Hence again applying the same convolution estimate we have
    \begin{equation}\label{eq:n=1_case_2}
        \begin{split}
            &\sum_{\eta_1 + \ldots + \eta_j \not=0} \sum_{\xi_1 + \ldots + \xi_j \not= 0} \lambda_{q+1}^{j(1-\epsilon_{q+1})\left(\frac{1}{k} - 1\right)} \left|\hat{\phi}_k\left(\lambda_{q+1}^{\epsilon_{q+1}-1}\xi_1\right) \cdots \hat{\phi}_k\left(\lambda_{q+1}^{\epsilon_{q+1}-1}\xi_j\right)\right| A_{q+1}^{\frac{j}{k}-1}\\
           &\times \left\Vert \hat{E}_q(t,\eta_1 - \lambda_{q+1}^{\epsilon_{q+1}} \xi_1)\right\Vert_{C_t^0}\delta(\eta_2 - \lambda_{q+1}^{\epsilon_{q+1}} \xi_2) \cdots \delta(\eta_j - \lambda_{q+1}^{\epsilon_{q+1}} \xi_j) |\eta_1 + \ldots + \eta_j|^{-s}\\
           &\lesssim  \sum_{\xi_1 + \ldots + \xi_j \not=0} \sum_{\eta_1 \not= \lambda_{q+1}^{\epsilon_{q+1}} \xi_1} \lambda_{q+1}^{j(1-\epsilon_{q+1})\left(\frac{1}{k}-1\right)-s\epsilon_{q+1}} A_{q+1}^{\frac{j}{k}-1} \left|\hat{\phi}_k\left(\lambda_{q+1}^{\epsilon_{q+1}-1}\xi_1\right) \cdots \hat{\phi}_k\left(\lambda_{q+1}^{\epsilon_{q+1}-1}\xi_j\right)\right|\\
           &\times \left\Vert \hat{E}_q(t,\eta_1 - \lambda_{q+1}^{\epsilon_{q+1}} \xi_1) \right\Vert_{C_t^0} |\xi_1 + \ldots + \xi_j|^{-s}\\
           &\lesssim \lambda_{q+1}^{(1-\epsilon_{q+1})\left(\frac{j}{k}-1\right) - s\epsilon_{q+1}} A_{q+1}^{\frac{j}{k}-1} \left(\sum_{\mu \not=0} \left\Vert \hat{E}_q(t,\mu) \right\Vert_{C_t^0}\right)\\
           &\lesssim \lambda_{q+1}^{(1-\epsilon_{q+1})\left(\frac{j}{k}-1\right) - s\epsilon_{q+1}} A_{q+1}^{\frac{j}{k}-1}
        \end{split}
    \end{equation}
    \textbf{Case 4:} $n_1 + n_2 + \ldots + n_j \geq 2$.
    Since $h_{q+1}\equiv 1$ on $\operatorname{supp}_t E_q$, we have
\[
    E_q^{n_i}h_{q+1}=E_q^{n_i}
    \qquad \text{whenever } n_i\geq 1,
\]
while for $n_i=0$,
\[
    \left(E_q^0h_{q+1}\right)^\wedge(t,\mu)
    =
    \delta_{\mu,0}h_{q+1}(t).
\]
Define
$$
C_q = \sum_{\xi \in \Z} \Vert \hat{E}_q(t,\xi) \Vert_{C_t^0} = \Vert E_q \Vert_{\A^0_{x,t}}.
$$
This quantity is finite due to the compact frequency support of $E_q$ discussed earlier and is clearly independent of $q+1$. Now using the algebra property of $\A^{0}_{x,t}$, we have
\[
    \|E_q^{n_i}h_{q+1}\|_{\dot{\mathbb{A}}_{x,t}^{-s}} < \Vert E_q^{n_i} \Vert_{\A^0_{x,t}}
    \lesssim C_q^{n_i}
\]
for every $n_i\geq0$.

Fix $n_1,\ldots,n_j\geq0$ with $n_1+\cdots+n_j\geq2$. We split according to whether $\xi_1+\cdots+\xi_j=0$, or not.

If $\xi_1+\cdots+\xi_j=0$, then
\[
    \eta_1+\cdots+\eta_j
    =
    \sum_{i=1}^j
    \left(\eta_i-\lambda_{q+1}^{\epsilon_{q+1}}\xi_i\right).
\]
After the change of variables $\mu_i=\eta_i-\lambda_{q+1}^{\epsilon_{q+1}}\xi_i$, we get
\[
\begin{aligned}
&\sum_{\eta_1+\cdots+\eta_j\neq0}
\prod_{i=1}^j
\left\|
\left(E_q^{n_i}h_{q+1}\right)^\wedge
\left(t,\eta_i-\lambda_{q+1}^{\epsilon_{q+1}}\xi_i\right)
\right\|_{C_t^0}
|\eta_1+\cdots+\eta_j|^{-s}
\\
&\qquad\le
\prod_{i=1}^j\sum_{\mu_i \in \Z}
\|(E_q^{n_i}h_{q+1})^\wedge(t,\mu_i)\|_{C_t^0}
\lesssim_q
C_q^{n_1+\cdots+n_j}.
\end{aligned}
\]
Therefore, using the same convolution estimate as in Case 2,
\begin{equation}\label{eq:page_29_est}
    \begin{split}
        &\sum_{\xi_1+\cdots+\xi_j=0}
\lambda_{q+1}^{j(1-\epsilon_{q+1})(\frac1k-1)}
A_{q+1}^{\frac jk-(n_1+\cdots+n_j)}
\left|\binom{1/k}{n_1}\cdots\binom{1/k}{n_j}\right|
\\
&\quad\times
\prod_{i=1}^j
\left|
\hat{\phi}_k
\left(\lambda_{q+1}^{\epsilon_{q+1}-1}\xi_i\right)
\right|
C_q^{n_1+\cdots+n_j}
\\
&\lesssim_{q,j,s}
\left|\binom{1/k}{n_1}\cdots\binom{1/k}{n_j}\right|
\lambda_{q+1}^{j(1-\epsilon_{q+1})(\frac1k-1)}
\lambda_{q+1}^{(j-1)(1-\epsilon_{q+1})}
A_{q+1}^{\frac jk-(n_1+\cdots+n_j)}
C_q^{n_1+\cdots+n_j}
\\
&=
\left|\binom{1/k}{n_1}\cdots\binom{1/k}{n_j}\right|
\lambda_{q+1}^{(1-\epsilon_{q+1})(\frac jk-1)}
A_{q+1}^{\frac jk}
\left(\frac{C_q}{A_{q+1}}\right)^{n_1+\cdots+n_j}.
    \end{split}
\end{equation}

Now suppose $\xi_1+\cdots+\xi_j\neq0$. Then, on the support of the frequency cutoff factors appearing in $\hat{w}_{q+1}$, each summand vanishes unless $\left|\eta_i-\lambda_{q+1}^{\epsilon_{q+1}}\xi_i \right| \lesssim \mu_{q+1}^{1/2}$ for $i=1,\cdots,j$. Hence,

\begin{equation}\label{eq:page29_est}
|\eta_1 +\cdots +\eta_j| \ge \lambda_{q+1}^{\epsilon_{q+1}}|\xi_1+ \cdots +\xi_j| - \sum_{i=1}^j \left(\eta_i-\lambda_{q+1}^{\epsilon_{q+1}}\xi_i\right) \ge \frac12 \lambda_{q+1}^{\epsilon_{q+1}}|\xi+\cdots+\xi_j|
\end{equation}
provided $\lambda_{q+1}$ is sufficiently large. Therefore
\[
|\eta_1 +\cdots +\eta_j|^{-s} \lesssim_s \lambda_{q+1}^{-s\epsilon_{q+1}}|\xi_1 + \cdots +\xi_j|^{-s}.
\]
Consequently,
\[
\begin{aligned}
    &\sum_{\eta_1,\cdots,\eta_j}\prod_{i=1}^j \Vert(E_q^{n_i}h_{q+1})^\wedge(t,\eta_i - \lambda_{q+1}^{\epsilon_{q+1}}\xi_i)\Vert_{C_t^0} |\eta_1 + \cdots + \eta_j|^{-s}\\
    &\qquad \lesssim_{j,s} \lambda_{q+1}^{-s\epsilon_{q+1}}|\xi_1 + \cdots + \xi_j|^{-s} \prod_{i=1}^j \sum_{\eta_i \in \Z} \Vert (E^{n_i}_q h_{q+1})^\wedge (t,\eta_i - \lambda_{q+1}^{\epsilon_{q+1}}\xi_i)\Vert_{C_t^0}\\
    &\qquad \lesssim_{q,j,s} \lambda_{q+1}^{-s\epsilon_{q+1}}|\xi_1 + \cdots + \xi_j|^{-s} C^{n_1+\cdots+n_j}_q.
\end{aligned}
\]
Using the same convolution estimate as in Case 1, we obtain
\begin{equation}\label{eq:page_29_est_3}
    \begin{split}
        &\sum_{\xi_1+\cdots+\xi_j\neq0}
\lambda_{q+1}^{j(1-\epsilon_{q+1})(\frac1k-1)}
A_{q+1}^{\frac jk-(n_1+\cdots+n_j)}
\left|\binom{1/k}{n_1}\cdots\binom{1/k}{n_j}\right|
\\
&\quad\times
\prod_{i=1}^j
\left|
\hat{\phi}_k
\left(\lambda_{q+1}^{\epsilon_{q+1}-1}\xi_i\right)
\right|
\lambda_{q+1}^{-s\epsilon_{q+1}}
|\xi_1+\cdots+\xi_j|^{-s}
C_q^{n_1+\cdots+n_j}
\\
&\lesssim_{q,j,s}
\left|\binom{1/k}{n_1}\cdots\binom{1/k}{n_j}\right|
\lambda_{q+1}^{j(1-\epsilon_{q+1})(\frac1k-1)-s\epsilon_{q+1}}
\lambda_{q+1}^{(j-1)(1-\epsilon_{q+1})}
A_{q+1}^{\frac jk-(n_1+\cdots+n_j)}
C_q^{n_1+\cdots+n_j}
\\
&=
\left|\binom{1/k}{n_1}\cdots\binom{1/k}{n_j}\right|
\lambda_{q+1}^{(1-\epsilon_{q+1})(\frac jk-1)-s\epsilon_{q+1}}
A_{q+1}^{\frac jk}
\left(\frac{C_q}{A_{q+1}}\right)^{n_1+\cdots+n_j}
\\
&\leq
\left|\binom{1/k}{n_1}\cdots\binom{1/k}{n_j}\right|
\lambda_{q+1}^{(1-\epsilon_{q+1})(\frac jk-1)}
A_{q+1}^{\frac jk}
\left(\frac{C_q}{A_{q+1}}\right)^{n_1+\cdots+n_j}.
    \end{split}
\end{equation}

Combining the two subcases, we obtain
\[
\begin{aligned}
&\sum_{\eta_1+\cdots+\eta_j\neq0}
\sum_{\substack{n_1,\ldots,n_j\geq0\\ n_1+\cdots+n_j\geq2}}
\sum_{\xi_1,\ldots,\xi_j\in\mathbb Z}
\lambda_{q+1}^{j(1-\epsilon_{q+1})(\frac1k-1)}
A_{q+1}^{\frac jk-(n_1+\cdots+n_j)}
\left|\binom{1/k}{n_1}\cdots\binom{1/k}{n_j}\right|
\\
&\quad\times
\prod_{i=1}^j
\left|
\hat{\phi}_k
\left(\lambda_{q+1}^{\epsilon_{q+1}-1}\xi_i\right)
\right|
\prod_{i=1}^j
\left\|
\left(E_q^{n_i}h_{q+1}\right)^\wedge
\left(t,\eta_i-\lambda_{q+1}^{\epsilon_{q+1}}\xi_i\right)
\right\|_{C_t^0}
|\eta_1+\cdots+\eta_j|^{-s}
\\
&\lesssim_{q,j,s}
\lambda_{q+1}^{(1-\epsilon_{q+1})(\frac jk-1)}
A_{q+1}^{\frac jk}
\sum_{\substack{n_1,\ldots,n_j\geq0\\ n_1+\cdots+n_j\geq2}}
\left|\binom{1/k}{n_1}\cdots\binom{1/k}{n_j}\right|
\left(\frac{C_q}{A_{q+1}}\right)^{n_1+\cdots+n_j}.
\end{aligned}
\]
Choosing $A_{q+1}$ sufficiently large so that
\begin{equation}\label{eq:A_q+1_cond}
    \frac{C_q}{A_{q+1}}\leq \frac{1}{2},
\end{equation}
the absolute convergence of the binomial series gives
\[
\sum_{\substack{n_1,\ldots,n_j\geq0\\ n_1+\cdots+n_j\geq2}}
\left|\binom{1/k}{n_1}\cdots\binom{1/k}{n_j}\right|
\left(\frac{C_q}{A_{q+1}}\right)^{n_1+\cdots+n_j}
\lesssim_{j,k}
\left(\frac{C_q}{A_{q+1}}\right)^2.
\]
Therefore
\[
\begin{aligned}
&\sum_{\eta_1+\cdots+\eta_j\neq0}
\sum_{\substack{n_1,\ldots,n_j\geq0\\ n_1+\cdots+n_j\geq2}}
\sum_{\xi_1,\ldots,\xi_j\in\mathbb Z}
\lambda_{q+1}^{j(1-\epsilon_{q+1})(\frac1k-1)}
A_{q+1}^{\frac jk-(n_1+\cdots+n_j)}
\left|\binom{1/k}{n_1}\cdots\binom{1/k}{n_j}\right|
\\
&\quad\times
\prod_{i=1}^j
\left|
\hat{\phi}_k
\left(\lambda_{q+1}^{\epsilon_{q+1}-1}\xi_i\right)
\right|
\prod_{i=1}^j
\left\|
\left(E_q^{n_i}h_{q+1}\right)^\wedge
\left(t,\eta_i-\lambda_{q+1}^{\epsilon_{q+1}}\xi_i\right)
\right\|_{C_t^0}
|\eta_1+\cdots+\eta_j|^{-s}
\\
&\lesssim_{q,j,s}
\lambda_{q+1}^{(1-\epsilon_{q+1})(\frac jk-1)}
A_{q+1}^{\frac jk}
\left(\frac{C_q}{A_{q+1}}\right)^2
\\
&\lesssim_{q,j,s}
\lambda_{q+1}^{(1-\epsilon_{q+1})(\frac jk-1)}
A_{q+1}^{\frac jk-2}.
\end{aligned}
\] 
Combining Cases 1-4, we obtain
\[
\begin{aligned}
AF(w_{q+1},j)
\lesssim_{q,j,s}\;&
\lambda_{q+1}^{(1-\epsilon_{q+1})(\frac jk-1)-s\epsilon_{q+1}}
A_{q+1}^{j/k}
\\
&+
\lambda_{q+1}^{(1-\epsilon_{q+1})(\frac jk-1)}
A_{q+1}^{\frac jk-1}
\|E_q\|_{\dot{\mathbb A}^{-s}_xC_t^0}
\\
&+
\lambda_{q+1}^{(1-\epsilon_{q+1})(\frac jk-1)-s\epsilon_{q+1}}
A_{q+1}^{\frac jk-1}
\\
&+
\lambda_{q+1}^{(1-\epsilon_{q+1})(\frac jk-1)}
A_{q+1}^{\frac jk-2}.
\end{aligned}
\]
Since $A_{q+1}\geq1$, $\lambda_{q+1}^{-s\epsilon_{q+1}}\leq1$, and
$\|E_q\|_{\mathbb{A}_{x,t}^{0}}$ is a $\lambda_{q+1}$-independent constant, we conclude
\[
AF(w_{q+1},j)
\lesssim_{q,j,s}
\lambda_{q+1}^{(1-\epsilon_{q+1})(\frac jk-1)}
A_{q+1}^{j/k}.
\]
\end{proof}

\end{document}
